Dans des conditions précises, le dihydrogène et le dichlore réagissent pour
donner du chlorure d'hydrogène de formule \ce{HCl} (composé gazeux).
\begin{enumerate}
    \item Écrire l'équation de la réaction.
    \item Les quantités initiales des réactifs sont~: $\qty{1.5}{\mol}$ de
          dihydrogène et $\qty{1.0}{\mol}$ de dichlore. Déterminer, en fonction
          de l'avancement $x$, les quantités de dihydrogène et de dichlore
          restantes et la quantité de chlorure d'hydrogène formé.
    \item Quel est le réactif limitant~?
    \item Quel est l'avancement maximal~?
    \item Quelle est la composition (en moles) du système à l'état final~?
\end{enumerate}

